\section{Finite certification and adversarial controls}

The executable certificate checks the determinant and local root law directly at
multiple scales and primes.  At $Q=25$, the determinant is $400$, the exceptional
prime $5$ has one root rather than two, the singular correction is exactly $4/3$, and
the unique edge is $(37,47)$.

A separate implementation scans every $8\le Q\le32768$.  It reproduces 568,308 edges
among 52,199,509 prime rows.  Table~\ref{tab:scan} records exact window statistics.
They are diagnostics, not evidence for the asymptotic theorem.

\begin{table}[ht]
\centering
\small
\begin{tabular}{rrrr}
\toprule
$Q$ window & edges & rows & aggregate $E/P$ \\
\midrule
$8$--$31$ & 3 & 115 & $3/115$ \\
$32$--$127$ & 32 & 1,472 & $1/46$ \\
$128$--$511$ & 338 & 19,171 & $338/19171$ \\
$512$--$2047$ & 3,681 & 255,321 & $409/28369$ \\
$2048$--$8191$ & 42,938 & 3,482,706 & $21469/1741353$ \\
$8192$--$32768$ & 521,316 & 48,440,724 & $43443/4036727$ \\
\bottomrule
\end{tabular}
\caption{Exact finite reproduction scan.}
\label{tab:scan}
\end{table}

The independent checker uses parameter-index enumeration rather than the producer's
prime-list enumeration.  Adversarial tests perturb the determinant by one, replace an
exceptional one-root prime by a generic two-root prime, and corrupt the Euler factor;
all three mutations are detected.  Normal and optimized Python modes are required to
produce identical output.
